\itemhead{习题4.2 · 第7题（教材p.114）}{XJ-B1-C04-HW-P114-Y0303}
\label{XJ-B1-C04-HW-P114-Y0303}
{\footnotesize 湘教版教材p.114 / PDF第121页 / 习题4.2 / 第7题}
\field{原题}{已知 $a,b,c,d$ 为不等于1的正数。图中为 $y=a^x$，$y=b^x$，$y=c^x$，$y=d^x$ 的图象，比较 $a,b,c,d$ 的大小：\blank。
\bankfig{C04-HW114}{0.55}}
\field{参考答案}{$b<a<1<d<c$，故填 $b<a<d<c$。在 $x=1$ 处读取函数值就是各底数；负半轴图象的上下关系相反。}
{\small \field{关联知识}{E01 指数函数的概念、图象与性质（教材帮p.150、151 / PDF159、160）；F04 函数三种表示法、图象与读数（教材帮p.85、86、93、94 / PDF94、95、102、103）}
\field{方法与易错点}{先查底数条件，利用已知点定底数，再据单调性及图象求值或范围。 易错：只要求底数非负；漏排1；将 $a^x+b$ 当标准指数函数且仍用原值域。}
\field{答案依据}{答案册相关题组：PDF第32页（册内页31）；按当前题库原题逐项重算，不按旧版题号机械对套。}}
